\begin{table}[!htbp]\centering\scriptsize
\caption{Les calendriers de la déduction immédiate comparés (Md€ ; part imposable 75~\%)}\label{tab:is_options}
\begin{tabularx}{\linewidth}{@{}>{\raggedright\arraybackslash}X>{\raggedleft\arraybackslash}p{1.45cm}>{\raggedleft\arraybackslash}p{1.45cm}>{\raggedleft\arraybackslash}p{1.45cm}>{\raggedleft\arraybackslash}p{1.45cm}>{\raggedleft\arraybackslash}p{1.45cm}@{}}
\toprule
Calendrier & Coût 2027 & Plus forte hausse annuelle & Coût actualisé par année d'achats & Dette maximale (cas central) & Extinction (central ; bornes) \\
\midrule
50~\% de 2027 à 2032, 100~\% en 2033 (ancienne note) & 7,0--7,8 & 7,9 & 1,7--2,6 & 11,2 & 2041 (2030 à 2051) \\
50~\% en 2027, +10 points par an, 100~\% en 2032 & 7,0--7,8 & 0,1 & 1,7--2,6 & 21,4 & 2042 (2036 à 2053) \\
100~\% sous 1 M€ dès 2027 ; au-delà, paliers 50 à 100~\% (2032) & 8,4--9,3 & 0,1 & 1,7--2,6 & 22,9 & 2042 (2036 à 2053) \\
100~\% sous 1 M€ par groupe, sans plus (AIA britannique seule) & 2,8--3,1 & 0,0 & 0,3--0,5 & aucune & --- \\
100~\% pour tous dès 2027 & 14,0--15,5 & 0,1 & 1,7--2,6 & 30,5 & 2043 (2036 à 2054) \\
50~\% pour tous, sans montée & 7,0--7,8 & 0,0 & 0,9--1,3 & 5,7 & 2033 (2030 à 2036) \\
Amortissement neutre (annuités revalorisées à 3,25~\%) & 0 & 0,4 & 1,1--1,6 & aucune & --- \\
\bottomrule
\end{tabularx}
\par\vspace{3pt}\parbox{\linewidth}{\raggedright\scriptsize Coût 2027 : régime actuel dégressif -- linéaire. Plus forte hausse annuelle : saut du coût d'une année sur l'autre (cas central) ; la falaise de 2033 en est le cas extrême. Coût actualisé à 5~\% (taux des entreprises), en régime permanent. Bornes de l'extinction : cas favorable (dégressif, ressources hautes) et défavorable (linéaire, ressources basses). Plafond : 1~M€ par groupe, couvrant 20~\% de l'investissement éligible (hypothèse). Amortissement neutre : plans actuels, annuités revalorisées au taux d'emprunt de l'État. Source : \texttt{is/calculs\_note\_is.py}.}
\end{table}
